\section{Технический проект}
\subsection{Архитектура программной системы}

Разрабатываемая платформа Easyappoint предназначена для онлайн-бронирования услуг частного мастера. Система обеспечивает регистрацию пользователей, настройку рабочей доступности мастера, создание пригласительных ссылок, просмотр свободных слотов, оформление записи клиента, синхронизацию с внешним календарем и отправку уведомлений.

Платформа построена по событийно-ориентированной микросервисной архитектуре. Данный подход позволяет отделить пользовательский интерфейс, основную бизнес-логику, хранение данных и доставку уведомлений. Основные операции бронирования выполняются синхронно через HTTP API, а уведомления и внешние интеграции обрабатываются асинхронно через событийный контур.

Архитектура системы включает следующие компоненты:
\begin{itemize}
\item клиентское web-приложение, реализующее интерфейсы мастера и клиента;
\item сервис данных \texttt{data-service}, отвечающий за аутентификацию, пользователей, приглашения, рабочую доступность, бронирования и интеграцию с Google Calendar;
\item реляционную базу данных PostgreSQL для хранения учетных записей, профилей, расписаний, записей, приглашений и служебных событий;
\item брокер сообщений Kafka для передачи доменных событий между сервисами;
\item сервис Telegram-уведомлений, который получает события записи и отправляет сообщения пользователям;
\item внешние сервисы Google Calendar API и Telegram Bot API.
\end{itemize}

Пользователь работает с системой через браузер. Клиентское приложение отправляет HTTP-запросы в сервис данных. Сервис данных проверяет права доступа, выполняет бизнес-логику и сохраняет результат в PostgreSQL. При создании, изменении или отмене записи сервис формирует доменное событие и помещает его в outbox-таблицу. Отдельный издатель отправляет подготовленные события в Kafka, после чего сервис Telegram-уведомлений получает событие и доставляет сообщение пользователю.

Общая схема архитектуры системы представлена на рисунке \ref{architecture:image}.

\begin{figure}[H]
\centering
\resizebox{0.98\linewidth}{!}{%
\begin{tikzpicture}[
component/.style={draw, rectangle, align=center, minimum width=3.4cm, minimum height=1.0cm, font=\small},
store/.style={draw, cylinder, shape border rotate=90, aspect=0.28, align=center, minimum width=2.9cm, minimum height=1.1cm, font=\small},
external/.style={draw, rectangle, dashed, align=center, minimum width=3.1cm, minimum height=0.95cm, font=\small},
relation/.style={-Latex, line width=0.65pt}
]
\node[component] (frontend) at (0,3.2) {Web-клиент\\React + TypeScript};
\node[component] (data) at (5.0,3.2) {Data Service\\Kotlin + Ktor};
\node[store] (db) at (10.0,3.2) {PostgreSQL};
\node[component] (kafka) at (5.0,0.8) {Kafka};
\node[component] (bot) at (10.0,0.8) {Telegram Bot\\Kotlin + Ktor};
\node[external] (google) at (10.0,5.6) {Google Calendar API};
\node[external] (telegram) at (14.3,0.8) {Telegram Bot API};

\draw[relation] (frontend) -- node[above, font=\scriptsize] {HTTP/JSON} (data);
\draw[relation] (data) -- node[above, font=\scriptsize] {SQL} (db);
\draw[relation] (data) -- node[right, font=\scriptsize] {OAuth, календарь} (google);
\draw[relation] (data) -- node[left, font=\scriptsize] {доменные события} (kafka);
\draw[relation] (kafka) -- node[above, font=\scriptsize] {booking.*} (bot);
\draw[relation] (bot) -- node[above, font=\scriptsize] {сообщения} (telegram);
\end{tikzpicture}%
}
\caption{Общая схема архитектуры системы}
\label{architecture:image}
\end{figure}

\subsection{Платформа онлайн-бронирования Easyappoint}

Платформа онлайн-бронирования Easyappoint является центральной частью разрабатываемой системы. Она объединяет работу мастера, клиента, внешнего календаря и канала уведомлений. В отличие от простого электронного журнала записей, система не только хранит бронирования, но и рассчитывает доступные интервалы с учетом рабочей доступности мастера, уже занятых слотов и данных внешнего календаря.

\subsubsection{Общая архитектура модуля}

Модуль построен по многоуровневой архитектуре и включает следующие логические слои:
\begin{itemize}
\item слой маршрутизации, принимающий HTTP-запросы, выполняющий первичную валидацию и передающий управление в соответствующий обработчик;
\item сервисный слой, в котором реализованы основные бизнес-правила: регистрация, создание приглашений, расчет слотов, оформление бронирования, отмена и изменение записи;
\item слой доступа к данным, выполняющий операции чтения и записи в PostgreSQL через jOOQ и пул соединений HikariCP;
\item интеграционный слой, отвечающий за работу с Google Calendar, Kafka и Telegram;
\item клиентский слой, реализованный в виде React-приложения с отдельными пользовательскими сценариями для мастера и клиента.
\end{itemize}

Такое разделение упрощает сопровождение системы: изменение интерфейса не требует переработки бизнес-логики, а изменение канала уведомлений не влияет на транзакционное сохранение записи.

\subsubsection{Обоснование технологического стека}

Серверная часть реализована на Kotlin с использованием Ktor. Kotlin выбран как статически типизированный язык, позволяющий компактно описывать бизнес-логику и снижать вероятность ошибок при работе с состояниями записей, ролями пользователей и временными интервалами. Ktor используется для построения HTTP API, настройки маршрутов, CORS, обработки ошибок и health-endpoint.

Клиентская часть реализована на React и TypeScript. TypeScript позволяет явно описывать структуры запросов и ответов, а React подходит для построения интерактивного web-интерфейса с отдельными кабинетами мастера и клиента. Для сборки клиентского приложения используется Vite.

Для хранения данных используется PostgreSQL. Схема базы данных управляется миграциями Flyway. Доступ к данным выполняется через jOOQ, что позволяет использовать типизированные SQL-запросы и явно контролировать транзакции. Для управления соединениями применяется HikariCP.

Событийный обмен реализован через Kafka. Он используется для передачи событий бронирования и уведомлений между сервисом данных и сервисом Telegram-уведомлений. Такой подход позволяет сохранить запись независимо от временной доступности внешнего канала связи.

\subsubsection{Реализация серверной части}

Серверная часть представлена сервисом \texttt{data-service}. Он является центральным компонентом системы и содержит основную бизнес-логику. Сервис запускается как самостоятельное Ktor-приложение и предоставляет REST API для клиентского web-приложения.

В сервисе данных выделены следующие функциональные модули:
\begin{itemize}
\item модуль аутентификации и авторизации, обеспечивающий регистрацию, вход, смену пароля и проверку JWT-токена;
\item модуль пользователей, отвечающий за данные текущего пользователя, профиль мастера и профиль клиента;
\item модуль приглашений, формирующий пригласительные ссылки для подключения клиентов к мастеру;
\item модуль рабочей доступности, рассчитывающий доступные интервалы мастера;
\item модуль бронирований, реализующий создание, изменение, отмену и просмотр записей;
\item модуль Google Calendar, отвечающий за подключение календаря, получение занятых интервалов и синхронизацию событий;
\item модуль уведомлений, управляющий привязкой Telegram-аккаунта;
\item модуль outbox, обеспечивающий надежную публикацию доменных событий в Kafka.
\end{itemize}

Обработка запроса выполняется последовательно. Сначала Ktor принимает HTTP-запрос и передает его в соответствующий route-обработчик. Далее выполняется проверка авторизации и роли пользователя. После этого управление передается в сервисный слой, где применяются бизнес-правила. Для операций с данными используются репозитории, работающие с PostgreSQL. Результат возвращается клиентскому приложению в формате JSON.

\subsubsection{Реализация клиентской части}

Клиентская часть реализована как React-приложение на TypeScript. Приложение предоставляет отдельные сценарии для мастера и клиента, а доступ к разделам ограничивается на уровне маршрутов и серверной авторизации.

В клиентском приложении выделены следующие группы экранов:
\begin{itemize}
\item публичные страницы входа и регистрации;
\item кабинет мастера с разделами обзора, календаря, приглашений, клиентов, записей и слотов;
\item кабинет клиента с просмотром доступности мастера и собственных записей;
\item общие настройки аккаунта: профиль, календарь, уведомления, безопасность и язык;
\item экран подключения по пригласительной ссылке.
\end{itemize}

Обмен данными с сервером вынесен в отдельные API-модули: \texttt{authApi}, \texttt{bookingApi}, \texttt{inviteApi}, \texttt{masterApi}, \texttt{googleApi}, \texttt{userApi}. Это разделяет сетевое взаимодействие и отображение данных на экране.

\subsubsection{Алгоритм расчета доступности}

Одним из основных компонентов платформы является модуль расчета доступных слотов. Он определяет интервалы, на которые клиент может записаться к мастеру. Расчет выполняется на основе часового пояса мастера, длительности слота, регулярных правил рабочей доступности, уже созданных активных записей и занятых интервалов внешнего календаря Google.

Алгоритм расчета включает следующие этапы.

Этап 1. Получение настроек расписания мастера.
Система загружает длительность слота, часовой пояс мастера и регулярные правила рабочей доступности. Эти данные определяют границы рабочего времени, в пределах которых могут быть сформированы потенциальные интервалы для записи.

Этап 2. Построение набора потенциальных слотов.
На основании выбранного клиентом периода и правил доступности система формирует предварительный набор временных интервалов. Каждый слот строится с учетом длительности записи и не выходит за пределы рабочего интервала мастера.

Этап 3. Получение активных бронирований.
Из базы данных выбираются существующие записи мастера за указанный период. При расчете учитываются только активные бронирования, поскольку отмененные записи не должны блокировать повторное использование временного интервала.

Этап 4. Получение занятых интервалов из Google Calendar.
Если у мастера подключен внешний календарь, система запрашивает занятость за выбранный период через Google Calendar API. Полученные интервалы используются как дополнительные ограничения доступности.

Этап 5. Исключение пересекающихся интервалов.
Активные бронирования и интервалы внешнего календаря объединяются в общий набор занятости. После этого из предварительного набора слотов исключаются все интервалы, которые пересекаются хотя бы с одним занятым промежутком.

Этап 6. Возврат доступных слотов клиентскому приложению.
После фильтрации система возвращает клиентскому приложению только свободные интервалы. Такой результат используется для отображения календаря записи и последующего создания бронирования.

Такой алгоритм позволяет мастеру задавать рабочую доступность заранее, а клиенту видеть только те интервалы, которые действительно могут быть забронированы.

\subsubsection{Управление бронированиями}

Создание бронирования выполняется с дополнительной серверной проверкой. Сервис убеждается, что пользователь имеет роль клиента, выбранный интервал корректен, приглашение действительно или клиент уже связан с мастером, а выбранный слот не пересекается с активной записью.

После успешной проверки запись сохраняется в базе данных, приглашение переводится в состояние принятого, при необходимости создается связь клиента с мастером, формируется outbox-событие и запускается синхронизация с Google Calendar.

Сценарий обработки бронирования представлен на рисунке \ref{booking-flow:image}.

\begin{figure}[H]
\centering
\resizebox{0.98\linewidth}{!}{%
\begin{tikzpicture}[
process/.style={draw, rectangle, align=center, minimum width=3.1cm, minimum height=0.95cm, font=\small},
store/.style={draw, cylinder, shape border rotate=90, aspect=0.28, align=center, minimum width=2.6cm, minimum height=1.0cm, font=\small},
relation/.style={-Latex, line width=0.65pt}
]
\node[process] (client) at (0,3.2) {Клиент\\выбирает слот};
\node[process] (api) at (4.2,3.2) {Booking API\\проверяет запрос};
\node[store] (db) at (8.4,3.2) {PostgreSQL};
\node[process] (google) at (4.2,1.1) {Google Calendar\\синхронизация};
\node[process] (outbox) at (8.4,1.1) {Outbox\\событие};
\node[process] (kafka) at (12.2,1.1) {Kafka};
\node[process] (bot) at (16.0,1.1) {Telegram Bot\\уведомляет};

\draw[relation] (client) -- (api);
\draw[relation] (api) -- node[above, font=\scriptsize] {запись} (db);
\draw[relation] (api) -- node[right, font=\scriptsize] {синхронизация} (google);
\draw[relation] (db) -- (outbox);
\draw[relation] (outbox) -- (kafka);
\draw[relation] (kafka) -- (bot);
\end{tikzpicture}%
}
\caption{Сценарий обработки бронирования}
\label{booking-flow:image}
\end{figure}

\subsubsection{Интеграция с Google Calendar}

Интеграция с Google Calendar используется для учета внешней занятости мастера и синхронизации созданных записей. Мастер подключает календарь через OAuth-механизм, после чего в базе данных сохраняются параметры подключения: идентификатор календаря, email аккаунта, access token, refresh token, область доступа и время истечения токена.

При расчете доступности сервис получает занятые интервалы из внешнего календаря и исключает их из доступных слотов. При создании или изменении записи сервис формирует событие календаря. При отмене записи соответствующее событие также обновляется или удаляется во внешнем календаре.

\subsubsection{Событийный обмен через Kafka}

Событийный обмен используется для отправки уведомлений о создании, изменении и отмене бронирований. В системе предусмотрены события \texttt{booking.created}, \texttt{booking.updated}, \texttt{booking.cancelled}, а также событие \texttt{notification.telegram.message} для отправки отдельных Telegram-сообщений.

Для надежной публикации событий применяется outbox-подход. При изменении состояния бронирования событие сначала сохраняется в таблицу \texttt{outbox\_events} в рамках транзакционной работы с базой данных. После завершения операции фоновый процесс публикует событие в Kafka. Если публикация не удалась, увеличивается счетчик попыток, сохраняется текст ошибки и назначается время следующей попытки.

\subsubsection{Уведомления Telegram}

Сервис Telegram-уведомлений подписан на Kafka-топики бронирований. При получении события он определяет тип операции, формирует текст уведомления по шаблону и отправляет сообщение через Telegram Bot API.

Пользователь подключает Telegram через одноразовый код. Для этого в базе данных создается запись в таблице \texttt{notification\_link\_tokens}. После подтверждения кода создается или обновляется канал доставки в таблице \texttt{notification\_destinations}. Такой подход позволяет хранить связь пользователя с Telegram отдельно от основной учетной записи.

\subsubsection{Управление сессиями и доступом}

Для управления пользовательской сессией используется JWT-токен. После успешного входа клиентское приложение получает токен доступа и прикладывает его к защищенным HTTP-запросам. Сервер проверяет подпись токена, идентификатор пользователя и роль.

В системе используются две основные роли: мастер и клиент. Мастер управляет профилем, расписанием, приглашениями, записями и интеграциями. Клиент просматривает доступные слоты мастера и управляет своими бронированиями. Проверка прав доступа выполняется на серверной стороне, поэтому клиентский интерфейс не является единственным механизмом защиты.

\subsubsection{Обработка ошибок и восстановление}

Для централизованной обработки ошибок используется механизм \texttt{StatusPages}. Ошибки валидации и нарушения бизнес-правил возвращаются пользователю в виде понятного сообщения, а непредвиденные исключения фиксируются в журнале и не раскрывают внутренние детали реализации.

Особое внимание уделено операциям бронирования и публикации событий. Создание записи выполняется в транзакции с проверкой пересечений и ограничений целостности. Если внешний календарь или Telegram временно недоступны, основная запись не теряется: данные сохраняются в PostgreSQL, а событие остается в outbox-таблице до следующей попытки публикации.

\subsubsection{Хранение данных, логирование и мониторинг}

Основные данные системы хранятся в PostgreSQL. Для управления схемой используются Flyway-миграции. При запуске сервиса данных миграции позволяют привести структуру базы данных к актуальному состоянию.

В сервисах предусмотрены эксплуатационные механизмы:
\begin{itemize}
\item health-endpoint для проверки доступности сервиса;
\item метрики Prometheus через Micrometer;
\item структурированное логирование;
\item correlation id для трассировки HTTP-запросов и Kafka-событий;
\item повторные попытки публикации outbox-событий.
\end{itemize}

\subsubsection{Заключение}

Разработанная платформа Easyappoint реализует полный цикл онлайн-бронирования: регистрацию пользователей, настройку рабочей доступности мастера, создание пригласительных ссылок, просмотр свободных слотов, оформление и изменение записей, синхронизацию с Google Calendar и доставку уведомлений через Telegram. Компоненты системы спроектированы с учетом расширяемости, транзакционной надежности и возможности дальнейшего развертывания в контейнерной среде.

\subsection{База данных и ее проектирование}

В качестве системы управления базами данных выбрана PostgreSQL. Эта СУБД обеспечивает транзакционность, ссылочную целостность, проверочные ограничения, индексы и возможность эффективной работы с временными интервалами. Для проекта онлайн-бронирования это важно, поскольку операции создания и изменения записей должны выполняться согласованно и не приводить к появлению некорректных данных.

Схема базы данных управляется через Flyway-миграции. Такой подход позволяет фиксировать изменения структуры таблиц в исходном коде, воспроизводимо разворачивать систему в разных средах и контролировать развитие модели данных.

\subsubsection{ER-модель данных}

На основе анализа предметной области выделены основные сущности системы: пользователи, мастера, клиенты, связи клиента с мастером, приглашения, правила рабочей доступности, бронирования, подключения Google Calendar и каналы уведомлений. Служебные таблицы токенов привязки уведомлений и outbox-событий относятся к физической реализации и подробно раскрываются в реляционной модели.

ER-модель базы данных представлена на рисунке \ref{er-model:image}.

\begin{figure}[H]
\centering
\includegraphics[width=\linewidth]{er-diagram.png}
\caption{ER-модель базы данных}
\label{er-model:image}
\end{figure}

Центральное место в модели занимают сущности \texttt{users}, \texttt{masters}, \texttt{clients} и \texttt{bookings}. Пользователь имеет одну роль, на основе которой создается профиль мастера или клиента. Мастер задает рабочую доступность и формирует приглашения, а клиент через приглашение получает возможность записаться на свободный слот. Таблица \texttt{bookings} связывает мастера, клиента и при необходимости приглашение. Линии на диаграмме показывают связи между сущностями и кратность этих связей. Сущности \texttt{google\_calendar\_connections} и \texttt{notification\_destinations} обеспечивают интеграции с календарем и доставку уведомлений.

\subsubsection{Нормализация}

Структура базы данных спроектирована в соответствии с требованиями третьей нормальной формы. Каждая таблица описывает одну логическую сущность, а связи между сущностями реализованы через внешние ключи.

Данные пользователя отделены от профилей мастера и клиента. Это позволяет хранить общие учетные сведения в таблице \texttt{users}, а роль-зависимые данные -- в таблицах \texttt{masters} и \texttt{clients}. Связь клиента с мастером вынесена в отдельную таблицу \texttt{master\_client\_links}, так как один мастер может обслуживать многих клиентов, а клиент может быть связан с разными мастерами.

Данные о расписании разделены на общие настройки \texttt{master\_schedule\_settings} и регулярные правила \texttt{master\_availability\_rules}. Это позволяет независимо хранить длительность слота и набор временных интервалов по дням недели. Сведения о внешнем календаре и уведомлениях вынесены в отдельные таблицы, так как они относятся к интеграциям и не должны перегружать основную таблицу пользователей.

\subsubsection{Реляционная модель данных}

На основе ER-модели построена реляционная модель базы данных. В ней каждая сущность представлена таблицей, а связи реализованы через первичные и внешние ключи. Реляционная модель представлена на рисунке \ref{relation-model:image}.

\begin{figure}[H]
\centering
\includegraphics[width=\linewidth,height=0.82\textheight]{relation-diagram.png}
\caption{Реляционная модель базы данных}
\label{relation-model:image}
\end{figure}

Направление стрелок на реляционной модели показывает, какая таблица содержит внешний ключ на родительскую таблицу. В таблицах используются первичные ключи типа \texttt{bigint}, внешние ключи для связи сущностей и ограничения целостности. Для ускорения запросов предусмотрены индексы по пользователям, мастерам, клиентам, статусам, времени записей и состоянию outbox-событий. Для поиска пересечений временных интервалов используется GIST-индекс по диапазону времени записи.

\subsubsection{Реализация базы данных}

Реализация базы данных выполнена с помощью SQL-миграций Flyway. Базовая схема создается первой миграцией, после чего последующие миграции расширяют модель данных: добавляют настройки расписания мастера, названия правил доступности, дополнительные поля профиля клиента, метаданные Telegram-уведомлений и уточненный уникальный индекс для повторной записи после отмены.

В описании таблиц ниже знак «*» в графе «Обяз.» означает, что атрибут не допускает значение \texttt{NULL}.

В таблицах \ref{db:users}--\ref{db:outbox} приведено подробное описание структуры таблиц базы данных.

{\small
\begin{xltabular}{\textwidth}{|p{1.7cm}|p{1.4cm}|p{3.5cm}|p{3.0cm}|X|}
\caption{Атрибуты таблицы «users» (Пользователи)\label{db:users}}\\ \hline
\centrow Тип ключа & \centrow Обяз. & \centrow Атрибут & \centrow Тип данных & \centrow Назначение \\ \hline
\endfirsthead
\continuecaption{Продолжение таблицы \ref{db:users}}
\centrow Тип ключа & \centrow Обяз. & \centrow Атрибут & \centrow Тип данных & \centrow Назначение \\ \hline
\finishhead
PK & * & \texttt{id} & \texttt{bigint} & Идентификатор пользователя. \\ \hline
UK &  & \texttt{telegram\_user\_id} & \texttt{bigint} & Идентификатор пользователя Telegram. \\ \hline
 &  & \texttt{username} & \texttt{varchar(128)} & Имя пользователя Telegram. \\ \hline
 &  & \texttt{first\_name} & \texttt{varchar(128)} & Имя пользователя. \\ \hline
 &  & \texttt{last\_name} & \texttt{varchar(128)} & Фамилия пользователя. \\ \hline
 & * & \texttt{role} & \texttt{varchar(16)} & Роль пользователя: \texttt{MASTER} или \texttt{CLIENT}. \\ \hline
UK & * & \texttt{email} & \texttt{varchar(255)} & Адрес электронной почты для входа. \\ \hline
 &  & \texttt{password\_hash} & \texttt{varchar(255)} & Хеш пароля. \\ \hline
 & * & \texttt{created\_at} & \texttt{timestamptz} & Дата создания записи. \\ \hline
 & * & \texttt{updated\_at} & \texttt{timestamptz} & Дата последнего обновления. \\ \hline
\end{xltabular}
}

{\small
\begin{xltabular}{\textwidth}{|p{1.7cm}|p{1.4cm}|p{3.5cm}|p{3.0cm}|X|}
\caption{Атрибуты таблицы «masters» (Мастера)\label{db:masters}}\\ \hline
\centrow Тип ключа & \centrow Обяз. & \centrow Атрибут & \centrow Тип данных & \centrow Назначение \\ \hline
\endfirsthead
\continuecaption{Продолжение таблицы \ref{db:masters}}
\centrow Тип ключа & \centrow Обяз. & \centrow Атрибут & \centrow Тип данных & \centrow Назначение \\ \hline
\finishhead
PK & * & \texttt{id} & \texttt{bigint} & Идентификатор профиля мастера. \\ \hline
FK, UK & * & \texttt{user\_id} & \texttt{bigint} & Ссылка на учетную запись пользователя. \\ \hline
 &  & \texttt{display\_name} & \texttt{varchar(255)} & Отображаемое имя мастера. \\ \hline
 &  & \texttt{description} & \texttt{text} & Описание мастера или оказываемых услуг. \\ \hline
 & * & \texttt{timezone} & \texttt{varchar(64)} & Часовой пояс мастера. \\ \hline
 & * & \texttt{created\_at} & \texttt{timestamptz} & Дата создания профиля. \\ \hline
 & * & \texttt{updated\_at} & \texttt{timestamptz} & Дата последнего обновления. \\ \hline
\end{xltabular}
}

{\small
\begin{xltabular}{\textwidth}{|p{1.7cm}|p{1.4cm}|p{3.5cm}|p{3.0cm}|X|}
\caption{Атрибуты таблицы «clients» (Клиенты)\label{db:clients}}\\ \hline
\centrow Тип ключа & \centrow Обяз. & \centrow Атрибут & \centrow Тип данных & \centrow Назначение \\ \hline
\endfirsthead
\continuecaption{Продолжение таблицы \ref{db:clients}}
\centrow Тип ключа & \centrow Обяз. & \centrow Атрибут & \centrow Тип данных & \centrow Назначение \\ \hline
\finishhead
PK & * & \texttt{id} & \texttt{bigint} & Идентификатор профиля клиента. \\ \hline
FK, UK & * & \texttt{user\_id} & \texttt{bigint} & Ссылка на учетную запись пользователя. \\ \hline
 &  & \texttt{display\_name} & \texttt{varchar(255)} & Отображаемое имя клиента. \\ \hline
 & * & \texttt{created\_at} & \texttt{timestamptz} & Дата создания профиля. \\ \hline
 & * & \texttt{updated\_at} & \texttt{timestamptz} & Дата последнего обновления. \\ \hline
\end{xltabular}
}

{\small
\begin{xltabular}{\textwidth}{|p{1.7cm}|p{1.4cm}|p{3.5cm}|p{3.0cm}|X|}
\caption{Атрибуты таблицы «master\_client\_links» (Связи клиентов с мастерами)\label{db:links}}\\ \hline
\centrow Тип ключа & \centrow Обяз. & \centrow Атрибут & \centrow Тип данных & \centrow Назначение \\ \hline
\endfirsthead
\continuecaption{Продолжение таблицы \ref{db:links}}
\centrow Тип ключа & \centrow Обяз. & \centrow Атрибут & \centrow Тип данных & \centrow Назначение \\ \hline
\finishhead
PK & * & \texttt{id} & \texttt{bigint} & Идентификатор связи. \\ \hline
FK, UK & * & \texttt{master\_id} & \texttt{bigint} & Ссылка на мастера, входит в уникальную пару. \\ \hline
FK, UK & * & \texttt{client\_id} & \texttt{bigint} & Ссылка на клиента, входит в уникальную пару. \\ \hline
 & * & \texttt{created\_at} & \texttt{timestamptz} & Дата создания связи. \\ \hline
\end{xltabular}
}

{\small
\begin{xltabular}{\textwidth}{|p{1.7cm}|p{1.4cm}|p{3.5cm}|p{3.0cm}|X|}
\caption{Атрибуты таблицы «invites» (Приглашения)\label{db:invites}}\\ \hline
\centrow Тип ключа & \centrow Обяз. & \centrow Атрибут & \centrow Тип данных & \centrow Назначение \\ \hline
\endfirsthead
\continuecaption{Продолжение таблицы \ref{db:invites}}
\centrow Тип ключа & \centrow Обяз. & \centrow Атрибут & \centrow Тип данных & \centrow Назначение \\ \hline
\finishhead
PK & * & \texttt{id} & \texttt{bigint} & Идентификатор приглашения. \\ \hline
FK & * & \texttt{master\_id} & \texttt{bigint} & Мастер, создавший приглашение. \\ \hline
FK &  & \texttt{client\_id} & \texttt{bigint} & Клиент, принявший приглашение. \\ \hline
UK & * & \texttt{token} & \texttt{varchar(128)} & Уникальный токен пригласительной ссылки. \\ \hline
 & * & \texttt{status} & \texttt{varchar(16)} & Статус: \texttt{ACTIVE}, \texttt{ACCEPTED}, \texttt{REVOKED}, \texttt{EXPIRED}. \\ \hline
 &  & \texttt{expires\_at} & \texttt{timestamptz} & Дата истечения приглашения. \\ \hline
 &  & \texttt{accepted\_at} & \texttt{timestamptz} & Дата принятия приглашения. \\ \hline
 &  & \texttt{revoked\_at} & \texttt{timestamptz} & Дата отзыва приглашения. \\ \hline
 & * & \texttt{created\_at} & \texttt{timestamptz} & Дата создания записи. \\ \hline
 & * & \texttt{updated\_at} & \texttt{timestamptz} & Дата последнего обновления. \\ \hline
\end{xltabular}
}

{\small
\begin{xltabular}{\textwidth}{|p{1.7cm}|p{1.4cm}|p{3.5cm}|p{3.0cm}|X|}
\caption{Атрибуты таблицы «master\_schedule\_settings» (Настройки расписания мастера)\label{db:schedule-settings}}\\ \hline
\centrow Тип ключа & \centrow Обяз. & \centrow Атрибут & \centrow Тип данных & \centrow Назначение \\ \hline
\endfirsthead
\continuecaption{Продолжение таблицы \ref{db:schedule-settings}}
\centrow Тип ключа & \centrow Обяз. & \centrow Атрибут & \centrow Тип данных & \centrow Назначение \\ \hline
\finishhead
PK & * & \texttt{id} & \texttt{bigint} & Идентификатор настроек. \\ \hline
FK, UK & * & \texttt{master\_id} & \texttt{bigint} & Мастер, которому принадлежат настройки. \\ \hline
 & * & \texttt{slot\_duration\_minutes} & \texttt{integer} & Длительность одного слота от 5 до 720 минут. \\ \hline
 & * & \texttt{created\_at} & \texttt{timestamptz} & Дата создания настроек. \\ \hline
 & * & \texttt{updated\_at} & \texttt{timestamptz} & Дата последнего обновления. \\ \hline
\end{xltabular}
}

{\small
\begin{xltabular}{\textwidth}{|p{1.7cm}|p{1.4cm}|p{3.5cm}|p{3.0cm}|X|}
\caption{Атрибуты таблицы «master\_availability\_rules» (Правила рабочей доступности)\label{db:availability-rules}}\\ \hline
\centrow Тип ключа & \centrow Обяз. & \centrow Атрибут & \centrow Тип данных & \centrow Назначение \\ \hline
\endfirsthead
\continuecaption{Продолжение таблицы \ref{db:availability-rules}}
\centrow Тип ключа & \centrow Обяз. & \centrow Атрибут & \centrow Тип данных & \centrow Назначение \\ \hline
\finishhead
PK & * & \texttt{id} & \texttt{bigint} & Идентификатор правила. \\ \hline
FK & * & \texttt{master\_id} & \texttt{bigint} & Мастер, которому принадлежит правило. \\ \hline
 & * & \texttt{day\_of\_week} & \texttt{integer} & День недели от 1 до 7. \\ \hline
 &  & \texttt{title} & \texttt{varchar(255)} & Название правила доступности. \\ \hline
 & * & \texttt{start\_time} & \texttt{time} & Время начала рабочего интервала. \\ \hline
 & * & \texttt{end\_time} & \texttt{time} & Время окончания рабочего интервала. \\ \hline
 & * & \texttt{created\_at} & \texttt{timestamptz} & Дата создания правила. \\ \hline
 & * & \texttt{updated\_at} & \texttt{timestamptz} & Дата последнего обновления. \\ \hline
\end{xltabular}
}

{\small
\begin{xltabular}{\textwidth}{|p{1.7cm}|p{1.4cm}|p{3.5cm}|p{3.0cm}|X|}
\caption{Атрибуты таблицы «bookings» (Бронирования)\label{db:bookings}}\\ \hline
\centrow Тип ключа & \centrow Обяз. & \centrow Атрибут & \centrow Тип данных & \centrow Назначение \\ \hline
\endfirsthead
\continuecaption{Продолжение таблицы \ref{db:bookings}}
\centrow Тип ключа & \centrow Обяз. & \centrow Атрибут & \centrow Тип данных & \centrow Назначение \\ \hline
\finishhead
PK & * & \texttt{id} & \texttt{bigint} & Идентификатор бронирования. \\ \hline
FK, UK & * & \texttt{master\_id} & \texttt{bigint} & Мастер, к которому оформлена запись. \\ \hline
FK, UK & * & \texttt{client\_id} & \texttt{bigint} & Клиент, оформивший запись. \\ \hline
FK &  & \texttt{invite\_id} & \texttt{bigint} & Приглашение, использованное при записи. \\ \hline
 & * & \texttt{status} & \texttt{varchar(16)} & Статус: \texttt{CREATED}, \texttt{CANCELLED}, \texttt{UPCOMING}, \texttt{COMPLETED}. \\ \hline
UK & * & \texttt{start\_at} & \texttt{timestamptz} & Начало интервала записи, входит в уникальный активный интервал. \\ \hline
UK & * & \texttt{end\_at} & \texttt{timestamptz} & Окончание интервала записи, входит в уникальный активный интервал. \\ \hline
UK &  & \texttt{google\_event\_id} & \texttt{varchar(512)} & Идентификатор события Google Calendar. \\ \hline
 &  & \texttt{cancelled\_at} & \texttt{timestamptz} & Дата отмены записи. \\ \hline
 &  & \texttt{cancellation\_reason} & \texttt{text} & Причина отмены записи. \\ \hline
 & * & \texttt{created\_at} & \texttt{timestamptz} & Дата создания записи. \\ \hline
 & * & \texttt{updated\_at} & \texttt{timestamptz} & Дата последнего обновления. \\ \hline
\end{xltabular}
}

{\small
\begin{xltabular}{\textwidth}{|p{1.7cm}|p{1.4cm}|p{3.5cm}|p{3.0cm}|X|}
\caption{Атрибуты таблицы «google\_calendar\_connections» (Подключения Google Calendar)\label{db:google}}\\ \hline
\centrow Тип ключа & \centrow Обяз. & \centrow Атрибут & \centrow Тип данных & \centrow Назначение \\ \hline
\endfirsthead
\continuecaption{Продолжение таблицы \ref{db:google}}
\centrow Тип ключа & \centrow Обяз. & \centrow Атрибут & \centrow Тип данных & \centrow Назначение \\ \hline
\finishhead
PK & * & \texttt{id} & \texttt{bigint} & Идентификатор подключения. \\ \hline
FK, UK & * & \texttt{master\_id} & \texttt{bigint} & Мастер, подключивший календарь. \\ \hline
 & * & \texttt{google\_calendar\_id} & \texttt{varchar(255)} & Идентификатор календаря Google. \\ \hline
 &  & \texttt{google\_account\_email} & \texttt{varchar(255)} & Email подключенного Google-аккаунта. \\ \hline
 &  & \texttt{access\_token} & \texttt{text} & Access token для обращения к API. \\ \hline
 &  & \texttt{refresh\_token} & \texttt{text} & Refresh token для обновления доступа. \\ \hline
 &  & \texttt{token\_type} & \texttt{varchar(64)} & Тип токена. \\ \hline
 &  & \texttt{scope} & \texttt{text} & Область доступа OAuth. \\ \hline
 &  & \texttt{access\_token\_expires\_at} & \texttt{timestamptz} & Время истечения access token. \\ \hline
 & * & \texttt{connected\_at} & \texttt{timestamptz} & Дата подключения календаря. \\ \hline
 &  & \texttt{disconnected\_at} & \texttt{timestamptz} & Дата отключения календаря. \\ \hline
 & * & \texttt{created\_at} & \texttt{timestamptz} & Дата создания записи. \\ \hline
 & * & \texttt{updated\_at} & \texttt{timestamptz} & Дата последнего обновления. \\ \hline
\end{xltabular}
}

{\small
\begin{xltabular}{\textwidth}{|p{1.7cm}|p{1.4cm}|p{3.5cm}|p{3.0cm}|X|}
\caption{Атрибуты таблицы «notification\_destinations» (Каналы уведомлений)\label{db:destinations}}\\ \hline
\centrow Тип ключа & \centrow Обяз. & \centrow Атрибут & \centrow Тип данных & \centrow Назначение \\ \hline
\endfirsthead
\continuecaption{Продолжение таблицы \ref{db:destinations}}
\centrow Тип ключа & \centrow Обяз. & \centrow Атрибут & \centrow Тип данных & \centrow Назначение \\ \hline
\finishhead
PK & * & \texttt{id} & \texttt{bigint} & Идентификатор канала уведомлений. \\ \hline
FK, UK & * & \texttt{user\_id} & \texttt{bigint} & Пользователь, которому принадлежит канал. \\ \hline
UK & * & \texttt{channel} & \texttt{varchar(32)} & Тип канала доставки, используется \texttt{TELEGRAM}. \\ \hline
 & * & \texttt{destination} & \texttt{varchar(255)} & Адрес доставки уведомления. \\ \hline
 &  & \texttt{destination\_label} & \texttt{varchar(255)} & Отображаемая подпись канала. \\ \hline
 &  & \texttt{verified\_at} & \texttt{timestamptz} & Дата подтверждения канала. \\ \hline
 & * & \texttt{enabled} & \texttt{boolean} & Признак включения канала. \\ \hline
 & * & \texttt{created\_at} & \texttt{timestamptz} & Дата создания записи. \\ \hline
 & * & \texttt{updated\_at} & \texttt{timestamptz} & Дата последнего обновления. \\ \hline
\end{xltabular}
}

{\small
\begin{xltabular}{\textwidth}{|p{1.7cm}|p{1.4cm}|p{3.5cm}|p{3.0cm}|X|}
\caption{Атрибуты таблицы «notification\_link\_tokens» (Коды привязки уведомлений)\label{db:notification-tokens}}\\ \hline
\centrow Тип ключа & \centrow Обяз. & \centrow Атрибут & \centrow Тип данных & \centrow Назначение \\ \hline
\endfirsthead
\continuecaption{Продолжение таблицы \ref{db:notification-tokens}}
\centrow Тип ключа & \centrow Обяз. & \centrow Атрибут & \centrow Тип данных & \centrow Назначение \\ \hline
\finishhead
PK & * & \texttt{id} & \texttt{bigint} & Идентификатор кода привязки. \\ \hline
FK & * & \texttt{user\_id} & \texttt{bigint} & Пользователь, которому выдан код. \\ \hline
 & * & \texttt{channel} & \texttt{varchar(32)} & Канал привязки, используется \texttt{TELEGRAM}. \\ \hline
UK & * & \texttt{code} & \texttt{varchar(64)} & Одноразовый код привязки. \\ \hline
 & * & \texttt{expires\_at} & \texttt{timestamptz} & Дата истечения кода. \\ \hline
 &  & \texttt{consumed\_at} & \texttt{timestamptz} & Дата использования кода. \\ \hline
 & * & \texttt{created\_at} & \texttt{timestamptz} & Дата создания записи. \\ \hline
\end{xltabular}
}

{\small
\begin{xltabular}{\textwidth}{|p{1.7cm}|p{1.4cm}|p{3.5cm}|p{3.0cm}|X|}
\caption{Атрибуты таблицы «outbox\_events» (Исходящие доменные события)\label{db:outbox}}\\ \hline
\centrow Тип ключа & \centrow Обяз. & \centrow Атрибут & \centrow Тип данных & \centrow Назначение \\ \hline
\endfirsthead
\continuecaption{Продолжение таблицы \ref{db:outbox}}
\centrow Тип ключа & \centrow Обяз. & \centrow Атрибут & \centrow Тип данных & \centrow Назначение \\ \hline
\finishhead
PK & * & \texttt{id} & \texttt{bigint} & Идентификатор события в outbox. \\ \hline
UK & * & \texttt{event\_id} & \texttt{uuid} & Уникальный идентификатор события. \\ \hline
 & * & \texttt{aggregate\_type} & \texttt{varchar(64)} & Тип агрегата, породившего событие. \\ \hline
 & * & \texttt{aggregate\_id} & \texttt{bigint} & Идентификатор агрегата. \\ \hline
 & * & \texttt{event\_type} & \texttt{varchar(128)} & Тип события, например \texttt{booking.created}. \\ \hline
 & * & \texttt{topic} & \texttt{varchar(255)} & Kafka-топик для публикации. \\ \hline
 &  & \texttt{message\_key} & \texttt{varchar(255)} & Ключ сообщения Kafka. \\ \hline
 & * & \texttt{payload} & \texttt{jsonb} & Полезная нагрузка события. \\ \hline
 & * & \texttt{headers} & \texttt{jsonb} & Заголовки события. \\ \hline
 & * & \texttt{status} & \texttt{varchar(16)} & Статус публикации: \texttt{PENDING}, \texttt{PUBLISHED}, \texttt{FAILED}. \\ \hline
 & * & \texttt{attempt\_count} & \texttt{integer} & Количество попыток публикации. \\ \hline
 & * & \texttt{next\_attempt\_at} & \texttt{timestamptz} & Время следующей попытки публикации. \\ \hline
 &  & \texttt{published\_at} & \texttt{timestamptz} & Дата успешной публикации. \\ \hline
 &  & \texttt{last\_error} & \texttt{text} & Текст последней ошибки публикации. \\ \hline
 & * & \texttt{created\_at} & \texttt{timestamptz} & Дата создания события. \\ \hline
\end{xltabular}
}

Таким образом, реализация базы данных Easyappoint включает основные предметные таблицы для пользователей, мастеров, клиентов, приглашений и бронирований, а также служебные таблицы для расписания, внешнего календаря, уведомлений и событийной публикации. Ограничения целостности и индексы позволяют поддерживать корректность данных и ускорять выполнение ключевых сценариев бронирования.
